\scanpage{102}
\uline{Azumaya rang 4}\hfill (II)\hfill Avril 86

Soit $A$ Alg. d'Azumaya rang $4$ sur \deleted{schéma $S$} \inserted{anneau \inserted{$k$ (disons)}}, et $B\subset A$ un sous-\deleted{Module loc. \unsure{ct.|f.} direct} \inserted{fibré vectoriel} de rang $3$. Alors $B^\perp$ est un sous-\deleted{\unsure{vectoriel} de} \inserted{fibré vectoriel} de rang $1$, soit $D$. On a corr. $1$-$1$ entre les $B$, et les $D$, car \unsure{les formes quadratiques|la forme quadratique} canonique $\operatorname{Tr}(u,v)$ sur $A$ est non dégénérée.

À vrai dire, il y a deux formes bil. sym. non dégénérées intéressantes sur $A$, savoir
\[
\varphi(u,v)=\operatorname{Tr}(uv)\qquad\text{et}\qquad\inserted{\psi(u,v)=}\varphi_{\det}(u,v)=\operatorname{Tr}u\operatorname{Tr}v-\operatorname{Tr}(uv),
\]
\deleted{donnant} \deleted{ayant des discriminants \unsure{qui ne}}.

Je dis que pour tt fibré vectoriel $B$ (de rang quelconque) qui contient $1$, \uline{les deux formes ont} $=$ \uline{discriminant sur} $B$. En effet, il suffit de le voir dans cas universel d'une grassmannienne sur $\operatorname{Spec}\mathbb{Z}$, ce qui nous permet de supposer $2$ inversible, donc
\[
A=k.1\oplus A^{\unsure{s|\leftrightarrow|\natural}}
\]
\margin{$A^s\overset{\text{déf}}{=}\{u\in A\mid\operatorname{Tr}u=0\}= \mathbb{1}^{\perp}$ (\uline{NB} L'orthogonal de $\mathbb{1}$ pour l'une ou l'autre forme c'est le $=$ !)}
Donc
\[
B=k.1\oplus V,\quad V\subset A^s
\]
% archive sheet number: 101
